% ============================================================================
%  DOCUMENTO PRINCIPAL — Actividad 7: Controlador de Dominio
%  ----------------------------------------------------------------------------
%  Aquí SOLO se integran los módulos. No se escribe contenido ni formato aquí.
%
%  Regla de oro del proyecto:
%    - formato/preambulo.sty  -> toda la configuración y paquetes (APA + Arial)
%    - formato/portada.tex    -> portada (datos institucionales)
%    - contenido/*.tex        -> el contenido real del informe
%    - anexos/*.tex           -> material complementario
%
%  Para compilar (desde latex/a7/):
%    pdflatex main && bibtex main && pdflatex main && pdflatex main
%  o directamente: make
% ============================================================================

\documentclass[12pt,a4paper]{article}

% --- Carga de la configuración de formato (APA + Arial) ---------------------
\usepackage{formato/preambulo}

% ============================================================================
%  DATOS DEL DOCUMENTO  (editar aquí)
% ============================================================================
\newcommand{\tituloDocumento}{Configuraci\'on Controlador de Dominio en Linux con Debian}
\newcommand{\equipo}{Grupo 1 Sudoers}
\newcommand{\integrantes}{%
  Joaquin Alessandro Felipez Rojas \par
  Nicolas Sebastian Reguerin Meneses \par
  David Willy Cruz Huanca}
\newcommand{\universidad}{Universidad Privada Domingo Savio}
\newcommand{\facultad}{Facultad de Ingenier\'ia}
\newcommand{\carrera}{Ingenier\'ia en Sistemas}
\newcommand{\asignatura}{Taller de Sistemas Operativos II}
\newcommand{\docente}{Jose Abraham Mita Huanca}
\newcommand{\ciudad}{Cochabamba--Bolivia}
\newcommand{\fechaEntrega}{23/09/2026}

\begin{document}

% --- Numeración: romano antes del Capítulo I, arábigo desde el Capítulo I ---
\pagenumbering{roman}

% --- Portada (no muestra número y no cuenta: el entorno titlepage parte de 1)
\input{formato/portada}

% --- Índices ----------------------------------------------------------------
\newpage
{\indicesCompactos\tableofcontents}
\newpage
{\indicesCompactos\listoftables}
\newpage
{\indicesCompactos\listoffigures}
\clearpage

% --- Cambio a numeración arábiga para el contenido principal -------------
\clearpage
\pagenumbering{arabic}

% ============================================================================
%  CONTENIDO  (un capítulo con su carátula y un archivo por sección)
% ============================================================================

% --- CAPÍTULO I --------------------------------------------------------------
\caratulaCapitulo{CAP\'ITULO I}{INTRODUCCI\'ON}
% ============================================================================
%  CAPÍTULO I: INTRODUCCIÓN
% ============================================================================

\section{Introducción}

En la gestión contemporánea de infraestructuras de tecnologías de la información (TI), la administración centralizada de identidades, credenciales y directivas de acceso constituye uno de los pilares fundamentales para garantizar la seguridad operativa, la escalabilidad y la eficiencia administrativa. Históricamente, las organizaciones enfrentaban serias dificultades derivadas de la gestión descentralizada de cuentas de usuario, donde cada estación de trabajo y servidor mantenía su propia base de datos de autenticación local. Este esquema distribuido no solo incrementaba exponencialmente la sobrecarga operativa del personal de sistemas al exigir la duplicación manual de usuarios y contraseñas en cada equipo, sino que también introducía vulnerabilidades críticas de seguridad, inconsistencias en los privilegios y una nula capacidad de auditoría unificada.

\subsection{El concepto de Controlador de Dominio en entornos Linux}

Frente a las limitaciones de los esquemas dispersos, surgió el modelo de Directorio Activo y de Controlador de Dominio (Domain Controller o DC), concebido como el núcleo autoritativo que almacena, estructura y valida la totalidad de identidades, relaciones de pertenencia a grupos y derechos de acceso dentro de un perímetro lógico denominado dominio \citep{msft_activedirectory}. En un Controlador de Dominio convergen de manera sinérgica tres servicios esenciales de red:

\begin{enumerate}[label=\alph*)]
    \item \textbf{Servicio de Directorio Ligero (LDAP):} Proporciona una base de datos jerárquica y optimizada para lecturas frecuentes, donde se organizan objetos tales como usuarios, grupos de seguridad, computadoras y unidades organizativas (OU), permitiendo consultas estructuradas y autorización granular \citep{rfc4511}.
    \item \textbf{Protocolo de autenticación Kerberos (v5):} Implementa un mecanismo criptográfico de boletos (\textit{tickets}) emitidos por un Centro de Distribución de Claves (KDC). Kerberos permite que usuarios y servicios demuestren mutuamente su identidad a través de redes no confiables sin transmitir contraseñas en texto claro, mitigando vectores de ataque por intercepción o suplantación \citep{rfc4120}.
    \item \textbf{Sistema de Nombres de Dominio (DNS):} Actúa como el sistema de localización y resolución fundamental de la red \citep{rfc1035}, extendido mediante registros de servicio (SRV) que posibilitan a las estaciones clientes descubrir dinámicamente qué servidores operan las funciones de autenticación LDAP y KDC dentro del espacio de nombres \citep{rfc2782}.
\end{enumerate}

Si bien estas arquitecturas fueron consolidadas y popularizadas en el ámbito corporativo por Microsoft Windows Server, el ecosistema de código abierto bajo GNU/Linux ha desarrollado soluciones maduras y robustas para asumir este rol neurálgico sin supeditar la infraestructura a licencias propietarias ni a ecosistemas cerrados.

\subsection{El rol crítico del Sistema de Nombres de Dominio (DNS)}

El Sistema de Nombres de Dominio (DNS) representa la columna vertebral de cualquier arquitectura distribuida. En las fases iniciales de diseño y laboratorio, la resolución del espacio de nombres interno de la organización (\texttt{sudoers.lan}) se estructuró mediante servidores independientes basados en BIND 9 (\textit{Berkeley Internet Name Domain}). Dicha configuración previa permitió establecer:

\begin{itemize}
    \item \textbf{Zonas directas e inversas:} Resolución de nombres a direcciones IPv4 mediante registros de host tipo A (\texttt{db.sudoers.lan}) y correspondencia inversa de direcciones IP a nombres canónicos mediante registros de puntero PTR en el dominio \texttt{in-addr.arpa} (\texttt{db.1.168.192}).
    \item \textbf{Mapeo de servicios departamentales:} Definición de registros A y CNAME para la totalidad de servicios corporativos de la red, tales como servidores web (\texttt{web}, \texttt{www}), correo electrónico (\texttt{mail} con su correspondiente registro MX de prioridad), almacenamiento compartido (\texttt{files}), servicios de impresión (\texttt{print}), monitoreo (\texttt{zabbix}) y administración de contenedores (\texttt{portainer}).
    \item \textbf{Vistas y reenviadores recursivos:} Implementación de vistas (\textit{views}) de BIND 9 para entornos con múltiples interfaces de red y reenviadores (\textit{forwarders}) hacia servidores públicos (\texttt{1.1.1.1}, \texttt{8.8.8.8}) para resolver peticiones fuera del dominio local.
\end{itemize}

No obstante, al dar el salto hacia una arquitectura de Directorio Activo, el rol del DNS trasciende la mera traducción estática de nombres a direcciones IP. En Active Directory, el DNS se convierte en un catálogo dinámico y crítico de localización de servicios: sin una resolución DNS impecable y especializada, ningún cliente puede encontrar el KDC para autenticarse ni consultar el catálogo global LDAP. Específicamente, el dominio depende de registros de servicio (SRV) bajo las ramas canónicas \texttt{\_msdcs}, \texttt{\_tcp} y \texttt{\_udp} (como \texttt{\_ldap.\_tcp.sudoers.lan} y \texttt{\_kerberos.\_tcp.sudoers.lan}). Por esta razón, el Controlador de Dominio asume directamente la autoridad sobre la zona \texttt{sudoers.lan}, garantizando la publicación y actualización continua de dichos registros vitales.

\subsection{Importancia e impacto del servicio Samba 4 Active Directory}

En este contexto de interoperabilidad y soberanía tecnológica, el proyecto Samba representa un hito fundamental en la historia del software libre. Durante décadas, Samba proveyó servicios de compartición de archivos e impresión bajo los protocolos SMB/CIFS y actuó como Controlador de Dominio Primario (PDC) al estilo clásico de Windows NT 4.0. No obstante, con la maduración de la arquitectura de Samba 4, la suite evolucionó para implementar de forma completa y nativa un Controlador de Dominio de Directorio Activo (\textit{Active Directory Domain Controller} o AD DC), plenamente compatible con los protocolos, esquemas y niveles funcionales de Microsoft Active Directory.

Para resolver la dependencia crítica de DNS en Active Directory, Samba 4 incorpora un backend DNS interno (\texttt{SAMBA\_INTERNAL}), o alternativamente un módulo dinámico para BIND 9 (BIND9\_DLZ). En la arquitectura adoptada, el servidor Samba AD gestiona directamente el servicio DNS interno en el puerto 53, asegurando:
\begin{itemize}
    \item \textbf{Gestión nativa de registros SRV y localizadores de dominio:} Administrando internamente los registros requeridos por Kerberos y LDAP sin requerir sincronizaciones manuales propensas a errores.
    \item \textbf{Continuidad en la resolución de servicios:} Permitiendo registrar y resolver todos los nombres de servicios de infraestructura heredados del laboratorio (\texttt{web}, \texttt{mail}, \texttt{print}, \texttt{files}) dentro de la misma zona autoritativa \texttt{sudoers.lan}.
    \item \textbf{Reenvío recursivo transparente:} Derivando consultas externas hacia servidores DNS globales para mantener acceso a Internet en las estaciones de la red.
\end{itemize}

La adopción de Samba 4 como AD DC sobre distribuciones empresariales como Debian GNU/Linux reviste una trascendencia estratégica por múltiples razones:
\begin{itemize}
    \item \textbf{Interoperabilidad multiplataforma transparente:} Al permitir que clientes heterogéneos (estaciones con GNU/Linux mediante SSSD o Winbind, y terminales con sistemas operativos Microsoft Windows) inicien sesión, validen credenciales y consuman recursos compartidos de manera idéntica y sin requerir agentes propietarios en los puestos de trabajo.
    \item \textbf{Eficiencia de costos y soberanía de licenciamiento:} Al eliminar la necesidad de adquirir costosas licencias de servidor y licencias de acceso de cliente (CALs), democratizando la implementación de arquitecturas empresariales seguras tanto en entornos corporativos como en ámbitos académicos.
    \item \textbf{Estabilidad, seguridad y control del sistema:} Al ejecutarse de manera nativa sobre un sistema operativo de grado empresarial como Debian GNU/Linux, se aprovechan las bondades del núcleo Linux en materia de rendimiento de red, aislamiento de privilegios, filtrado granular por cortafuegos (\textit{nftables}) y automatización de tareas mediante herramientas de línea de comandos e infraestructuras reproducibles.
\end{itemize}

\subsection{Propósito de la actividad}

El presente trabajo, desarrollado en el marco de la asignatura \textit{Taller de Sistemas Operativos II}, tiene como propósito fundamental investigar, diseñar, implementar y documentar la puesta en marcha de un Controlador de Dominio corporativo basado en Samba 4 AD DC desplegado de forma nativa sobre una distribución Debian GNU/Linux. 

A través de esta práctica integral, se busca consolidar las competencias técnicas requeridas para desplegar el árbol de identidades de una organización real, articular la transición desde un servidor DNS estático hacia el motor de resolución DNS interno autoritativo de Active Directory para el dominio local \texttt{sudoers.lan}, aplicar políticas de seguridad en la gestión de contraseñas, estructurar recursos compartidos de red con permisos segmentados por grupos departamentales y verificar experimentalmente el inicio de sesión centralizado y la unión de estaciones clientes al dominio. El informe detalla de manera rigurosa cada decisión arquitectónica adoptada, contrastando los fundamentos teóricos con las evidencias prácticas obtenidas durante el proceso de configuración.

% --- CAPÍTULO II -------------------------------------------------------------
\caratulaCapitulo{CAP\'ITULO II}{OBJETIVOS, ALCANCE Y L\'IMITES}
% ============================================================================
%  CAPÍTULO II: OBJETIVOS, ALCANCE Y LÍMITES
%  Contenido pendiente de redacción.
% ============================================================================

\section{Objetivos del trabajo}

[Escribir aqu\'i los objetivos de la Actividad 7.]
% ============================================================================
%  CAPÍTULO II: ALCANCE Y LÍMITES
%  Contenido pendiente de redacción.
% ============================================================================

\section{Alcance y L\'imites}

[Escribir aqu\'i el alcance y los l\'imites de la Actividad 7.]

% --- CAPÍTULO III ------------------------------------------------------------
\caratulaCapitulo{CAP\'ITULO III}{DESARROLLO DEL TRABAJO}
% ============================================================================
%  CAPÍTULO III: DESARROLLO DEL TRABAJO
% ============================================================================

\section{Desarrollo del trabajo}

Este capítulo documenta, paso a paso, la configuración de un Controlador de
Dominio de Directorio Activo sobre Debian GNU/Linux 12. Cada bloque describe
desde la preparación del sistema operativo base hasta las pruebas de
autenticación y de unión de clientes al dominio, incluyendo la justificación
técnica de cada una de las decisiones tomadas.

% ============================================================================
\section{Preparación del sistema}

La primera etapa consistió en dejar el servidor listo para asumir el rol de
Controlador de Dominio. Un AD DC es el punto de referencia de identidad de
toda la red, por lo que su nombre y su dirección IP deben ser estables y
conocidos de antemano \citep{msft_activedirectory}.

\subsection{Identidad del host}

A los clientes de la red les llega el DC por \textit{hostname}; si este
cambiara, todos los equipos que ya lo conocen perderían la referencia. Por
ello se fijó el nombre \texttt{dc1} y el nombre completamente cualificado
(FQDN) \texttt{dc1.sudoers.lan}:

\begin{lstlisting}[style=terminal]
$ hostnamectl set-hostname dc1
$ hostnamectl status | head -3
  Static hostname: dc1
  Icon name: computer-laptop
  Operating System: Debian GNU/Linux 12 (bookworm)
\end{lstlisting}

\capturaAPA{fig:hostname}{Identidad de host \texttt{dc1} configurada en el servidor}{captura-hostname.png}{Captura de pantalla mostrando la salida de \texttt{hostnamectl status} con el nombre de host \texttt{dc1}.}

\subsection{Configuración de la IP fija}

Un DC no puede tener una dirección IP variable: Kerberos y el DNS lo
referencian constantemente por \texttt{192.168.0.2}. Además, Samba exige que
el \textit{hostname} resuelva a la IP de la red local y \textbf{nunca} a
\texttt{127.0.0.1} — si resolviera a \textit{loopback}, el DC se ``hablaría
a sí mismo'' y los clientes jamás lo alcanzarían. Por eso se garantizó que
\texttt{dc1} resuelva a \texttt{192.168.0.2/24} en la interfaz cableada
\texttt{eno1} y se registró en \texttt{/etc/hosts}:

\begin{lstlisting}[style=terminal]
$ ip -4 a show eno1
2: eno1: <BROADCAST,MULTICAST,UP,LOWER_UP> ... state UP
    inet 192.168.0.2/24 brd 192.168.0.255 scope global eno1

$ grep -E 'dc1|sudoers' /etc/hosts
192.168.0.2   dc1.sudoers.lan   dc1
\end{lstlisting}

\capturaAPA{fig:ip-fija}{Dirección IP fija \texttt{192.168.0.2/24} en la interfaz \texttt{eno1}}{captura-ip-fija.png}{Captura de pantalla mostrando la salida de \texttt{ip -4 a show eno1} con la IP \texttt{192.168.0.2/24} y la entrada de \texttt{/etc/hosts} que asocia \texttt{dc1.sudoers.lan} a esa dirección.}

\subsection{Instalación de paquetes y servicios del host}

Los paquetes del dominio (Samba 4, cliente Kerberos y utilidades DNS) se
instalaron de forma \textbf{nativa} sobre el host Debian, nunca dentro de un
contenedor: el DC de Samba es un servicio frágil cuya virtualización
contenerizada está descartada por decisión de diseño. La instalación se
realiza con el administrador APT:

\begin{lstlisting}[style=terminal]
$ apt update
$ apt install -y samba smbclient winbind krb5-user \
    libpam-winbind libnss-winbind dnsutils
\end{lstlisting}

Junto con Samba se habilitaron los demás roles del servidor ``todo-en-uno''
que dan sustento físico al dominio: el firewall \texttt{nftables}, el
servidor de hora \texttt{chrony} (imprescindible para Kerberos, que exige un
desfase de reloj menor a 5 minutos) y Docker para los servicios
contenedorizados del resto del proyecto:

\begin{lstlisting}[style=terminal]
$ systemctl is-enabled samba-ad-dc nftables chrony docker
samba-ad-dc.service    enabled
nftables.service       enabled
chrony.service         enabled
docker.service         enabled
\end{lstlisting}

\capturaAPA{fig:servicios-host}{Servicios del host habilitados para el arranque}{captura-servicios-host.png}{Captura de pantalla mostrando la salida de \texttt{systemctl is-enabled} con los servicios \texttt{samba-ad-dc}, \texttt{nftables}, \texttt{chrony} y \texttt{docker} en estado \texttt{enabled}.}

% ============================================================================
\section{Instalación de Samba 4 y promoción del servidor como DC}

Samba 4 integra en un único daemon (\texttt{samba-ad-dc}) todos los
protocolos que forman Active Directory: LDAP (directorio) \citep{rfc4511},
Kerberos (autenticación) \citep{rfc4120}, SMB/CIFS (archivos) y el DNS
propio \citep{sambateam2024}. El servidor se promueve a Controlador de
Dominio mediante el aprovisionamiento del dominio con \texttt{samba-tool}:

\begin{lstlisting}[style=terminal]
$ samba-tool domain provision \
    --use-rfc2307 --realm=SUDOERS.LAN \
    --domain=SUDOERS --server-role=dc \
    --adminpass='AdminPass!2026'
\end{lstlisting}

El aprovisionamiento crea el directorio LDAP, los relatios Kerberos, el
archivo \texttt{/etc/samba/smb.conf} y la maqueta DNS de la zona
\texttt{sudoers.lan}. Tras la promoción, se configura el cliente Kerberos
(\texttt{/etc/krb5.conf}) apuntando al KDC local (\texttt{dc1.sudoers.lan})
y se deja \texttt{samba-ad-dc} como el único demonio del dominio activo,
inhabilitando los heredados \texttt{smbd}, \texttt{nmbd} y \texttt{winbind}
para evitar colisiones de sockets \citep{sambateam2024}.

La promoción correcta se verifica con dos consultas de
\texttt{samba-tool}. La primera muestra el nivel funcional, es decir el
``modo de compatibilidad'' con el que opera el dominio (cuanto más alto,
más funciones activas ofrece):

\begin{lstlisting}[style=terminal]
$ samba-tool domain level show
Domain: SUDOERS
Domain level: 2008 R2
Forest level: 2008 R2
\end{lstlisting}

La segunda pregunta al propio DC \emph{dónde está} y \emph{a qué dominio
sirve}; la línea \texttt{Server role: active directory domain controller}
es la prueba de que la promoción como Controlador de Dominio funcionó:

\begin{lstlisting}[style=terminal]
$ samba-tool domain info 127.0.0.1
Domain:           SUDOERS
Domain SID:       S-1-5-21-...
Forest:           SUDOERS.LAN
Number of users:  13
Server role:      active directory domain controller
\end{lstlisting}

\capturaAPA{fig:domain-info}{Datos del dominio y del Controlador de Dominio}{captura-domain-info.png}{Captura de pantalla mostrando la salida de \texttt{samba-tool domain info 127.0.0.1} con el \textit{realm} \texttt{SUDOERS.LAN}, el nombre NetBIOS \texttt{SUDOERS} y el rol \texttt{active directory domain controller}.}

Finalmente, se comprueba que los puertos del dominio estén escuchando:
\textbf{53} (DNS), \textbf{88} (Kerberos), \textbf{389} (LDAP) y \textbf{445}
(SMB). Es la prueba física de red de que el DC está operativo:

\begin{lstlisting}[style=terminal]
$ ss -ltn | grep -E ':53|:88|:389|:445'
LISTEN  0   128  *:88    *:*   samba  (kerberos)
LISTEN  0   128  *:389   *:*   samba  (ldap)
LISTEN  0   128  *:445   *:*   samba  (smb)
LISTEN  0   128  *:53    *:*   samba  (dns)
\end{lstlisting}

\capturaAPA{fig:puertos-ad}{Puertos del dominio escuchando: 53, 88, 389 y 445}{captura-puertos-ad.png}{Captura de pantalla mostrando la salida de \texttt{ss -ltn} con los puertos 53, 88, 389 y 445 a la escucha.}

% ============================================================================
\section{Configuración del DNS interno}

Un dominio de Active Directory se apoya por completo en el DNS: los clientes
lo consultan para descubrir en qué servidor se encuentra cada servicio de
dominio \citep{rfc1035}. Por ello el DC de Samba asumió la autoridad de la
zona \texttt{sudoers.lan} con el backend \textbf{SAMBA\_INTERNAL} (gestionado
automáticamente por el propio DC, sin competir con el BIND9 que se retiró del
laboratorio).

\subsection{Registro A del Controlador de Dominio}

El registro \texttt{A} es la tabla nombre-$\rightarrow$IP del dominio: toda
consulta empieza pidiendo \texttt{dc1.sudoers.lan}:

\begin{lstlisting}[style=terminal]
$ dig dc1.sudoers.lan
;; ANSWER SECTION:
dc1.sudoers.lan. 900 IN A 192.168.0.2
\end{lstlisting}

\capturaAPA{fig:dig-dc}{Registro A del DC: \texttt{dc1.sudoers.lan} $\rightarrow$ \texttt{192.168.0.2}}{captura-dig-dc.png}{Captura de pantalla mostrando la salida de \texttt{dig dc1.sudoers.lan} con la respuesta \texttt{192.168.0.2}.}

\subsection{Registros SRV de descubrimiento}

Los registros SRV son las ``guías de servicios'' del dominio \citep{rfc2782}:
le dicen a cualquier cliente dónde está el LDAP, el Kerberos o el DC. Sin
ella las estaciones no encuentran el dominio. Su presencia confirma que el
descubrimiento automático funciona sin configuración manual:

\begin{lstlisting}[style=terminal]
$ dig _ldap._tcp.dc._msdcs.sudoers.lan SRV
;; ANSWER SECTION:
_ldap._tcp.dc._msdcs.sudoers.lan. 900 IN SRV 0 100 389 dc1.sudoers.lan.
\end{lstlisting}

\capturaAPA{fig:dig-srv}{Registros SRV de descubrimiento del dominio}{captura-dig-srv.png}{Captura de pantalla mostrando la salida de \texttt{dig \_ldap.\_tcp.dc.\_msdcs.sudoers.lan SRV} con el registro SRV apuntando al puerto 389 de \texttt{dc1.sudoers.lan}.}

\subsection{Zona completa del dominio}

Además de los registros de infraestructura, la zona guarda las direcciones
``lindas'' de los servicios y de los miembros del proyecto
(\texttt{joaquin.sudoers.lan}, \texttt{web.sudoers.lan}, \texttt{print.},
\texttt{mail.}, etc.). Todos resuelven de forma centralizada en el DC:

\begin{lstlisting}[style=terminal]
$ kinit joaquin@SUDOERS.LAN           # ticket para autenticar la consulta
$ samba-tool dns query dc1.sudoers.lan sudoers.lan @ ALL -k yes
Name=, Records=8, Children=0
 DC@SUDOERS.LAN: type=A, value=192.168.0.2
 web@SUDOERS.LAN: type=A, value=192.168.0.2
 joaquin@SUDOERS.LAN: type=A, value=192.168.0.2
...
\end{lstlisting}

\capturaAPA{fig:zona-dns}{Zona DNS completa de \texttt{sudoers.lan} gestionada por el DC}{captura-zona-dns.png}{Captura de pantalla mostrando la salida de \texttt{samba-tool dns query 127.0.0.1 sudoers.lan @ ALL} con los registros A de la zona.}

\subsection{Reenviador hacia Internet}

Para que los clientes resuelvan también nombres externos, el DC funciona con
un \textit{forwarder} hacia el router de la red. De esta manera una única
respuesta DNS del DC cubre tanto los nombres internos (zona
\texttt{sudoers.lan}) como los externos de Internet \citep{iscbind9}.

% ============================================================================
\section{Creación de usuarios y grupos}

Los permisos se otorgan a \textbf{grupos}, no a personas sueltas: de esa
forma ``quién puede qué'' se administra en un solo lugar y los cambios se
propagan a toda la organización. Los grupos y usuarios se crearon con el
script idempotente \texttt{dc/add-users-groups.sh} (puede ejecutarse varias
veces sin duplicar objetos) o manualmente con \texttt{samba-tool}.

\begin{lstlisting}[style=terminal]
$ samba-tool group list | grep -E 'admins|sistemas|oficina|contabilidad'
admins
sistemas
contabilidad
oficina

$ samba-tool user list
joaquin
nicolas
david
grupo2
grupo3
...
\end{lstlisting}

\capturaAPA{fig:grupos}{Grupos de seguridad del dominio}{captura-grupos.png}{Captura de pantalla mostrando la salida de \texttt{samba-tool group list} filtrada por los grupos \texttt{admins}, \texttt{sistemas} y \texttt{oficina}.}

\capturaAPA{fig:usuarios}{Usuarios registrados en el dominio}{captura-usuarios.png}{Captura de pantalla mostrando la salida de \texttt{samba-tool user list} con los usuarios administradores y de prueba.}

Los usuarios son objetos dentro del árbol LDAP, no cuentas sueltas en cada
máquina. Cada objeto tiene un \texttt{sAMAccountName} (el login corto), un
principal de usuario UPN (\texttt{joaquin@SUDOERS.LAN}) y pertenece a
\textit{n} grupos, heredando la suma de sus permisos:

\begin{lstlisting}[style=terminal]
$ samba-tool user show joaquin
dn: CN=joaquin,CN=Users,DC=SUDOERS,DC=LAN
sAMAccountName: joaquin
userPrincipalName: joaquin@SUDOERS.LAN
memberOf: CN=admins,CN=Users,DC=SUDOERS,DC=LAN
          CN=sistemas,CN=Users,DC=SUDOERS,DC=LAN
\end{lstlisting}

\capturaAPA{fig:user-show}{Detalle de un usuario y su pertenencia a grupos}{captura-user-show.png}{Captura de pantalla mostrando la salida de \texttt{samba-tool user show joaquin} con el atributo \texttt{memberOf} enumerando los grupos del usuario.}

La tabla \ref{tab:decisiones} resume las decisiones técnicas más relevantes
de las etapas anteriores y su justificación.

\begin{table}[H]
\centering
\caption{Decisiones técnicas y justificación}
\label{tab:decisiones}
\contenidoConFuente{%
  \begin{tablaAPA}
    \begin{tabular}{@{}p{5cm}p{8.5cm}@{}}
    \toprule
    \textbf{Decisión} & \textbf{Justificación} \\
    \midrule
    DC nativo (fuera de Docker) & Samba 4 como DC es frágil en contenedores; se despliega directamente sobre Debian \citep{sambateam2024} \\
    IP fija \texttt{192.168.0.2} / hostname \texttt{dc1} & Kerberos y DNS referencian constantemente al DC; un cambio de IP o nombre rompe el dominio \\
    \texttt{samba-tool domain provision} & Método oficial de promoción a DC; genera configuración, directorio LDAP y maqueta DNS \\
    DNS \textbf{SAMBA\_INTERNAL} autoritativo & El AD exige que el DC sea quien sirva la zona \texttt{sudoers.lan} y los registros SRV \\
    Permisos por grupos AD & Administración centralizada de ``quién puede qué''; herencia automática de permisos \\
    \texttt{samba-ad-dc} como único demonio & Evita colisiones de puertos con \texttt{smbd}, \texttt{nmbd} y \texttt{winbind} heredados \\
    \bottomrule
    \end{tabular}
  \end{tablaAPA}%
}{Elaboración propia.}
\end{table}

% ============================================================================
\section{Pruebas de autenticación}

El protocolo Kerberos resuelve el problema de ``demostrar quién soy sin
enviar la contraseña por la red'' \citep{rfc4120}. La respuesta son
\textbf{tickets}: el KDC (el propio DC) emite un ticket de concesión (TGT) al
autenticar al usuario, y ese ticket viaja firmado por el DC, de modo que
nadie puede falsificar una identidad.

La prueba en el servidor se realiza con las herramientas del paquete
\texttt{krb5-user}:

\begin{lstlisting}[style=terminal]
$ kinit juan@SUDOERS.LAN
Password for juan@SUDOERS.LAN:
$ klist
Ticket cache: FILE:/tmp/krb5cc_1000
Default principal: juan@SUDOERS.LAN
Valid starting     Expires            Service principal
09/20/26 10:00:00  09/20/26 20:00:00  krbtgt/SUDOERS.LAN@SUDOERS.LAN
$ kdestroy
\end{lstlisting}

\capturaAPA{fig:kinit-klist}{Solicitud y verificación de ticket Kerberos (TGT)}{captura-kinit-klist.png}{Captura de pantalla mostrando la salida de \texttt{kinit} y \texttt{klist} sobre \texttt{dc1}, con el ticket emitido por el KDC \texttt{SUDOERS.LAN} y su vigencia.}

\subsection{Autenticación desde un cliente real}

La misma cuenta y el mismo ``carnet'' funcionan desde cualquier equipo de la
red, porque el emisor de tickets siempre es el DC. Esto es el
*single sign-on*: el cliente se identifica una vez y accede a los servicios
del dominio \citep{msft_activedirectory}.

\begin{lstlisting}[style=terminal]
$ sudo kinit juan@SUDOERS.LAN && klist   # desde el cliente unido
Default principal: juan@SUDOERS.LAN
Valid starting     Expires            Service principal
09/20/26 10:05:00  09/20/26 20:05:00  krbtgt/SUDOERS.LAN@SUDOERS.LAN
\end{lstlisting}

\capturaAPA{fig:kinit-cliente}{Autenticación Kerberos desde un equipo cliente}{captura-kinit-cliente.png}{Captura de pantalla ejecutando \texttt{kinit} y \texttt{klist} en el equipo cliente, mostrando el ticket obtenido contra el DC.}

% ============================================================================
\section{Unión de clientes al dominio}

Unirse al dominio es más que configurar el DNS \citep{msft_activedirectory}:
\texttt{realm join} crea una \textbf{computer account} en el AD, convirtiendo
al equipo en un objeto del dominio con su propio secreto de máquina. Desde
entonces el cliente usa el AD como fuente de identidad (LDAP para leer,
Kerberos para autenticar) sin copiar usuarios locales.

\subsection{Unión del cliente Debian}

El cliente se unió al dominio ejecutando \texttt{realm join --user=administrator
SUDOERS.LAN} desde el equipo, autenticándose con la cuenta de administración
del dominio. La operación fue exitosa y quedó registrada la \emph{computer
account} correspondiente, como se verifica en el servidor más adelante.

\subsection{Verificación desde el cliente}

Los usuarios del AD aparecen como si fueran locales gracias a SSSD:
\texttt{getent} los consulta en vivo contra el DC, sin duplicar cuentas en el
equipo.

\begin{lstlisting}[style=terminal]
$ getent passwd joaquin david nicolas
joaquin:*:10001:10001:...:/home/joaquin:/bin/bash
david:*:10002:10002:...:/home/david:/bin/bash
nicolas:*:10003:10003:...:/home/nicolas:/bin/bash

$ getent group 'admins'
admins:*:20001:joaquin,david,nicolas
\end{lstlisting}

\capturaAPA{fig:getent}{Usuarios del dominio visibles desde el cliente con \texttt{getent}}{captura-getent.png}{Captura de pantalla mostrando la salida de \texttt{getent passwd} y \texttt{getent group} con los usuarios del AD resueltos en el cliente.}

\subsection{Verificación desde el servidor}

Del lado del server se comprueba que la notebook haya quedado registrada como
equipo del dominio:

\begin{lstlisting}[style=terminal]
$ samba-tool computer list
DESKTOP-NBK1$
...
\end{lstlisting}

\capturaAPA{fig:computer-list}{Estaciones registradas en el dominio}{captura-computer-list.png}{Captura de pantalla mostrando la salida de \texttt{samba-tool computer list} con el equipo del cliente unido al AD.}

\subsection{Acceso a recursos de archivo}

El DC sirve shares de forma nativa (sin contenedor) con permisos definidos
por grupos AD en el servidor: \texttt{departamentos}, \texttt{homes} (cada
usuario ve su propio directorio) y \texttt{respaldo}. El acceso se verifica
desde el servidor con \texttt{smbclient} y desde el cliente mediante un
montaje CIFS o el gestor de archivos:

\begin{lstlisting}[style=terminal]
$ smbclient -L dc1 --user=joaquin          # en el server
	Sharename       Type      Comment
	---------       ----      -------
	departamentos   Disk      Recursos por departamentos
	homes           Disk      Directorios personales
	respaldo        Disk
\end{lstlisting}

\capturaAPA{fig:smbclient}{Listado de los recursos compartidos (SMB) del DC}{captura-smbclient.png}{Captura de pantalla mostrando la salida de \texttt{smbclient -L dc1 --user=joaquin} con los shares \texttt{departamentos}, \texttt{homes} y \texttt{respaldo}.}

% ============================================================================
\section{DHCP y DNS en las estaciones de la red}

Para que una estación recién conectada quede ``lista para el dominio'' sin
intervención manual, la red cuenta con el servidor DHCP Kea (en contenedor,
con \texttt{network\_mode: host} porque DHCP usa broadcast y no atraviesa un
bridge de Docker). Kea entrega junto con la IP las \textbf{opciones} de red:
DNS \texttt{192.168.0.2}, dominio \texttt{sudoers.lan}, gateway
\texttt{192.168.0.1} y NTP \texttt{192.168.0.2}. El DHCP del router está
desactivado para evitar el doble servidor DHCP.

\subsection{Verificación del servicio DHCP}

\begin{lstlisting}[style=terminal]
$ sudo ss -ulnp | grep ':67 '
UNCONN  0  0   192.168.0.2:67  *:*  users:(("kea-dhcp4",pid=...))

$ sudo cat /var/lib/kea/kea-leases4.csv   # concesiones en vivo
address,hostname,state,valid_lifetime,...
192.168.0.100,notebook1,assigned,...
\end{lstlisting}

\capturaAPA{fig:kea-ss}{Servidor DHCP Kea escuchando en \texttt{192.168.0.2:67}}{captura-kea-ss.png}{Captura de pantalla mostrando la salida de \texttt{ss -ulnp} con Kea a la escucha en el puerto 67/udp y las concesiones activas.}

\subsection{Prueba de obtención de IP desde un cliente}

Un cliente se desconecta y vuelve a tomar IP, quedando con la configuración
completa entregada por Kea:

\begin{lstlisting}[style=terminal]
$ sudo dhclient -v eno2
DHCPDISCOVER on eno2 to 255.255.255.255 port 67
DHCPOFFER from 192.168.0.2
DHCPACK from 192.168.0.2 (xid=0x...)
$ ip -4 a
inet 192.168.0.101/24 ...   # DNS = 192.168.0.2, dominio = sudoers.lan
$ dig dc1.sudoers.lan      # resuelve con el DNS entregado por DHCP
\end{lstlisting}

\capturaAPA{fig:dhclient}{Cliente tomando IP del pool de Kea y resolviendo el DC}{captura-dhclient.png}{Captura de pantalla mostrando \texttt{dhclient -v} obteniendo una IP del pool dinámico y la posterior consulta \texttt{dig dc1.sudoers.lan} resuelta con el DNS entregado por DHCP.}

% --- CONCLUSIONES ---------------------------------------------------------
\newpage
\encabezadoCapitulo{CONCLUSIONES}{}
\addcontentsline{toc}{part}{CONCLUSIONES}
\begingroup
\setcounter{capitulo}{4}
\setcounter{section}{0}
\renewcommand{\thesection}{4.}
\renewcommand{\theHsection}{4.}
\renewcommand{\thesubsection}{4.\arabic{subsection}}
\renewcommand{\theHsubsection}{4.\arabic{subsection}}
\renewcommand{\thesubsubsection}{4.\arabic{subsection}.\arabic{subsubsection}}
\renewcommand{\theHsubsubsection}{4.\arabic{subsection}.\arabic{subsubsection}}
% ============================================================================
%  CAPÍTULO IV: CONCLUSIONES
% ============================================================================

\section{Conclusiones}

La implementación del Controlador de Dominio con Samba 4 sobre Debian
permitió comprender de manera integral cómo funciona la identidad centralizada
en una organización y qué servicios la hacen posible: el directorio LDAP, la
autenticación Kerberos, el DNS del dominio y el protocolo SMB/CIFS operando
como un único daemon. A continuación se presentan las conclusiones más
relevantes.

\subsection{Resultados obtenidos}

El Controlador de Dominio fue configurado y verificado exitosamente con las
siguientes características:

\begin{itemize}
  \item Samba 4 desplegado de forma \textbf{nativa} sobre Debian 12, con
        el daemon \texttt{samba-ad-dc} como único servicio del dominio
        (LDAP + Kerberos + SMB + DNS en un solo proceso).
  \item Dominio \texttt{SUDOERS.LAN} aprovisionado con
        \texttt{samba-tool domain provision}, nivel funcional 2008 R2 y rol
        verificado como \texttt{active directory domain controller}.
  \item DNS interno autoritativo con backend \textbf{SAMBA\_INTERNAL}: la
        zona \texttt{sudoers.lan} y los registros SRV de descubrimiento
        responden correctamente, permitiendo a cualquier cliente encontrar el
        dominio sin configuración manual.
  \item Usuarios (\texttt{joaquin}, \texttt{nicolas}, \texttt{david} y los
        de prueba \texttt{grupo2} a \texttt{grupo9}) y grupos de seguridad
        (\texttt{admins}, \texttt{sistemas}, \texttt{oficina}, entre otros)
        creados de forma idempotente con \texttt{dc/add-users-groups.sh}.
  \item Autenticación Kerberos verificada desde el propio DC y desde equipos
        clientes (\texttt{kinit}/\texttt{klist}), confirmando que el ticket
        se emite siempre contra el KDC del dominio.
  \item Unión de clientes Debian al dominio mediante \texttt{realm join},
        verificada por ambos lados: \texttt{getent} en el cliente y
        \texttt{samba-tool computer list} en el servidor.
  \item Acceso a los recursos compartidos nativos (\texttt{departamentos},
        \texttt{homes}, \texttt{respaldo}) con permisos definidos por grupos
        AD.
  \item El DHCP Kea entrega en el pool dinámico las opciones que dejan a cada
        estación lista para el dominio (DNS, dominio, gateway y NTP) sin
        intervención manual.
\end{itemize}

\subsection{Dificultades encontradas}

Durante la configuración se presentaron las siguientes dificultades:

\begin{itemize}
  \item \textbf{Coincidencia DNS \texttt{/etc/hosts}}: el DC debe resolver su
        propio nombre a la IP de la red y nunca a \texttt{127.0.0.1}; si
        resuelve a \textit{loopback}, el dominio queda ``hablando solo'' y los
        clientes no lo alcanzan.
  \item \textbf{Orden de arranque de servicios}: los servicios del host (DC,
        Docker, DNS, DHCP) y las interfaces de red deben iniciarse en orden;
        fue necesario retrasar la red y asegurar \texttt{network-online.target}
        para que Kea y los contenedores no arrancaran antes que la interfaz.
  \item \textbf{Kerberos y sincronización de reloj}: la autenticación falla si
        el reloj del cliente se desvía más de 5 minutos del KDC; se requirió
        chrony en el servidor y configurar NTP en las estaciones.
  \item \textbf{Broadcast y DHCP}: el DHCP de Kea exige
        \texttt{network\_mode: host}, porque BOOTP/DHCP usa broadcast y no
        puede mapearse a través de un bridge de Docker.
\end{itemize}

\subsection{Importancia del Controlador de Dominio en la administración de sistemas}

Un Controlador de Dominio es la base sobre la que se apoya la administración
de identidades de una organización:

\begin{itemize}
  \item \textbf{Centralización de la identidad}: una única cuenta por usuario
        sirve para autenticarse en cualquier equipo y acceder a los recursos,
        aplicando el concepto de *single sign-on* \citep{msft_activedirectory}.
  \item \textbf{Administración de permisos por grupos}: los permisos se
        otorgan a grupos y no a personas, reduciendo la superficie de error y
        centralizando el ``quién puede qué''.
  \item \textbf{Descubrimiento automático}: gracias a los registros SRV del
        DNS del DC, un equipo encuentra los servicios del dominio sin
        configuración manual \citep{rfc2782}.
  \item \textbf{Seguridad por diseño}: Kerberos elimina el envío de
        contraseñas por la red, ya que toda autenticación se basa en tickets
        firmados por el KDC \citep{rfc4120}.
  \item \textbf{Base para el resto de la infraestructura}: el DC sostiene a
        los demás servicios del proyecto (web, correo, compartición de
        archivos y DHCP), y una resolución DNS o un reloj mal configurados
        impactan en todos ellos.
\end{itemize}

\subsection{Aprendizajes adquiridos}

El trabajo permitió desarrollar competencias prácticas en administración de
servicios de red y de identidad:

\begin{itemize}
  \item Se comprendió la arquitectura de Active Directory y la función de cada
        protocolo que integra Samba 4 (LDAP, Kerberos, SMB y DNS) funcionando
        como un único daemon.
  \item Se experimentó en primera persona que el DNS es la base del dominio:
        sin registros A y SRV correctos ningún cliente logra encontrar ni
        autenticarse contra el DC \citep{rfc1035}.
  \item Se valoró la administración nativa de servicios en Debian y el uso de
        herramientas propias como \texttt{samba-tool}, lo que brinda un
        entendimiento profundo de cada componente del sistema.
  \item Se aprendió a integrar el DC con el resto del ecosistema del
        proyecto (Docker, nftables, chrony, Kea) respetando decisiones de
        diseño no negociables, como que el DC nunca se conteneriza y que sus
        puertos no se re-publican en contenedores.
  \item Se aplicó el concepto de infraestructura como código: la promoción
        del DC, los usuarios y los grupos quedaron definidos en scripts y
        archivos de configuración versionados (\texttt{dc/provision.sh},
        \texttt{dc/add-users-groups.sh}), permitiendo reproducir el entorno de
        forma exacta y predecible (como en el lab \texttt{TSO-II}).
  \item Se comprendió la equivalencia de la solución con un AD de Microsoft
        (igual REINO, Nivel Funcional, \texttt{getent}/\texttt{ldapsearch}
        como \texttt{dsquery}) tanto para administradores como para clientes
        Windows unidos al dominio.
\end{itemize}

Como trabajo a futuro, queda pendiente la integración de los clientes
Windows de contabilidad y marketing, la implementación de políticas de equipo
(GPO) equivalentes a las de Windows Server y la puesta en producción del
respaldo y monitoreo del DC nativo.
\endgroup

% --- Referencias (APA) ------------------------------------------------------
\newpage
\encabezadoCapitulo{REFERENCIAS}{}
\addcontentsline{toc}{part}{REFERENCIAS}
\bibliographystyle{apacite}
\nocite{*}
\bibliography{formato/referencias}

% --- Anexos (material complementario, tras las referencias) ------------------
% \newpage
% \encabezadoCapitulo{ANEXOS}{}
% \pdfbookmark[0]{ANEXOS}{anexos}%
% \addcontentsline{toc}{part}{ANEXOS}
% \appendix
% % ============================================================================
%  ANEXO A: [TÍTULO DEL ANEXO]
%  Plantilla: copiar para cada anexo (anexo-a, anexo-b, ...).
% ============================================================================

\section{Anexo A: [T\'itulo del anexo]}

[Escribir el contenido del anexo aqu\'i.]

% Ejemplo de tabla:
% \begin{table}[H]
%   \centering
%   \caption{T\'itulo de la tabla}
%   \contenidoConFuente{%
%     \begin{tablaAPA}
%       \begin{tabular}{@{}p{3.6cm}p{8.2cm}p{3.2cm}@{}}
%         \toprule
%         \textbf{Columna 1} & \textbf{Columna 2} & \textbf{Columna 3} \\
%         \midrule
%         Dato 1 & Dato 2 & Dato 3 \\
%         \bottomrule
%       \end{tabular}
%     \end{tablaAPA}%
%   }{Elaboraci\'on propia.}
% \end{table}

\end{document}